\documentclass[12pt]{article}

\usepackage[a4paper, margin=2.5cm]{geometry}
\usepackage{amsmath}
\usepackage{amssymb}
\usepackage{amsthm}
\usepackage[colorlinks=true, linkcolor=blue, citecolor=blue, urlcolor=blue]{hyperref}
\usepackage{booktabs}
\usepackage{natbib}

\newtheorem{theorem}{Theorem}
\newtheorem{proposition}[theorem]{Proposition}
\newtheorem{corollary}[theorem]{Corollary}
\newtheorem{definition}[theorem]{Definition}
\newtheorem{remark}[theorem]{Remark}
\newtheorem{example}[theorem]{Example}
\newtheorem{lemma}[theorem]{Lemma}

\newcommand{\Ecal}{\mathcal{E}}
\newcommand{\Gcal}{\mathcal{G}}
\newcommand{\Fcal}{\mathcal{F}}
\newcommand{\HW}{H_{\mathrm{W}}}
\newcommand{\Rst}{R_{\mathrm{st}}}

\title{A Structural Bridge Hypothesis\\Between Quantum and Gravitational Event-Order Regimes}
\author{%
  Lee Smart\\[4pt]
  \textit{Institute of Vibrational Field Dynamics}\\[2pt]
  \texttt{contact@vibrationalfielddynamics.org}\\[2pt]
  \texttt{@vfd\_org}%
}
\date{April 2026}

\begin{document}

\maketitle

\noindent\textbf{Status:} Preprint (not peer-reviewed)

\begin{abstract}
The VFD programme has developed two structural tracks from related geometric substrates: a \textit{quantum track} (Papers~XV--XXII) in which a Schr\"odinger equation is recovered at equilibrium and at formal leading order near equilibrium under the Nelson pairing construction (Papers~XVIII, XXI), and a \textit{gravitational track} (Papers~XXIII--XXVII) constructing discrete scaffold analogues of temporal ordering (XXIII, conditional), observer frames (XXIV, under finite-poset hypotheses), Hasse-graph metric separation (XXV), graph-curvature indicators (XXVI), and a graph-Poisson source/potential model (XXVII).

The present paper argues, at \emph{structural-correspondence level}, that these two tracks can be jointly viewed through a shared set of name-level objects (event/configuration set, closure functional, graph Laplacian). The paper does not derive Einstein equations from quantum mechanics or vice versa. It does not establish that the quantum-track and gravitational-track Laplacians act on the same space, nor that the closure functional used by Paper~XXII coincides with the admissibility criterion used by Papers~XXIII--XXVII. Paper~XXI explicitly notes that Route~A and Route~B of the quantum track use different kinetic operators on different spaces; Papers~XXIII--XXVII's geometric content is discrete scaffold, not a derived gravitational dynamics.

A \emph{pattern-correspondence table} is presented listing proposed role-parallels across the two tracks (with an explicit ``Status'' column flagging which rows are formal analogies on different spaces versus genuine mathematical identifications; all rows are analogies). A \emph{combinatorial regime indicator} $\eta$ (ordering fraction) is proposed as a bookkeeping statistic on event posets; $\eta$ is not a derived physical control parameter and no dynamical crossover between the two tracks is proved.

This paper proposes a structural bridge hypothesis, not a unification theorem. It identifies a pattern within the VFD programme and flags the open identification questions that would be required to upgrade the pattern to a substrate-level identification theorem.
\end{abstract}

%==================================================================
\section{Introduction}\label{sec:intro}
%==================================================================

Two apparently distinct physical regimes --- quantum behaviour and gravitational structure --- have resisted unification within the standard framework for nearly a century. This paper does not resolve that problem. It argues, at \emph{structural-correspondence level}, that within the VFD closure programme the mathematical objects named in the quantum track (Papers~XV--XXII) and the gravitational track (Papers~XXIII--XXVII) carry parallel role-labels that invite a common substrate hypothesis; it does not establish such a common substrate as a theorem.

\subsection*{The two tracks}

The VFD programme has constructed two tracks independently:

\begin{enumerate}
    \item \textbf{Quantum track} (Papers~XV--XXII). Starting from the stochastic closure dynamics of Paper~XIV~\citep{SmartXIV}, this track constructs a complexified-action ansatz~\citep{SmartXV}, proves a dissipative obstruction~\citep{SmartXVII}, resolves it via the Nelson pairing construction~\citep{SmartXVIII}, records norm-related residual properties~\citep{SmartXIX,SmartXX}, synthesises the construction (with Route~A on closure-state space and the conditional Route~B on the 600-cell vertex graph)~\citep{SmartXXI}, and connects to the 600-cell spectral structure~\citep{SmartXXII}. The \emph{regime-restricted result}: Schr\"odinger recovery is exact at Paper~XIV equilibrium and formal at leading order near equilibrium; global linear Schr\"odinger recovery is explicitly \emph{not} established in the cited papers.

    \item \textbf{Gravitational track} (Papers~XXIII--XXVII). From the dual 600-cell phase-overlap structure, this track constructs discrete scaffold analogues of temporal ordering (XXIII, basepoint-dependent strict order, conditional on incomparable events), observer-frame kinematics (XXIV, under finite-poset hypotheses, not Lorentz-derived), Hasse-graph metric separation on connected components (XXV), graph-curvature indicators on toy models (XXVI), and a graph-Poisson source-potential model (XXVII). The \emph{regime-restricted result}: each paper supplies one discrete scaffold piece; no paper derives gravitational dynamics proper. Paper~XXVII explicitly disclaims an Einstein equation, stress-energy tensor, continuum limit, and geodesic-from-$u$ law.
\end{enumerate}

\subsection*{Central observation}

Both tracks were constructed from related 600-cell-derived structures, using objects with overlapping names but track-specific definitions. This paper makes the name-level overlap explicit and asks: is this an accident of presentation, or does it reflect a deeper structural identification?

\subsection*{Scope}

This paper maps candidate name-level correspondences across the two tracks, proposes a combinatorial regime indicator, and formulates a substrate-bridge hypothesis. It does not derive the Einstein equations from quantum mechanics, does not produce a theory of quantum gravity, and does not claim to resolve the measurement problem. It proposes a structural bridge hypothesis between two tracks that were built from related 600-cell-derived structures.

%==================================================================
\section{Shared Bookkeeping Vocabulary}\label{sec:substrate}
%==================================================================

Both tracks can be discussed using a common bookkeeping vocabulary over related 600-cell-derived structures, comprising three name-level objects defined with track-specific content in the preceding papers.

\subsection{The dual 600-cell system}

Paper~XXII~\citep{SmartXXII} records a $240$-to-$120$ two-to-one projection (in the counting/projection sense; the underlying theorem-grade construction is in the substrate paper~\citep{CascadeAlgebraicSubstrate}). Up to a fixed/degenerate-pair caveat discussed in the substrate source, the 240 roots of $E_8$ project $2$-to-$1$ onto the 120 vertices of the 600-cell under the Galois involution of the icosian construction; this is an involutive counting pairing, not a topological double cover. The quantum track uses the \textit{spectral structure} of the 600-cell graph (eigenvalues, representations, closure functional). The gravitational track uses the \textit{overlap structure} of the dual system (admissible events from relative phase rotation).

\begin{remark}[Related structures within the dual 600-cell construction]
The two tracks access related structures within the dual 600-cell construction: spectral data (quantum: 600-cell graph eigenvalues/representations) and event-order data (gravitational: admissible overlaps from two phase-rotated copies, Paper~XXIII). No same-object identification across the two tracks is proved; Paper~XXIII explicitly works with two phase-rotated 600-cell copies rather than identifying them with the quantum-track 600-cell spectral structure.
\end{remark}

\subsection{The closure functional}

The quantum track uses a specific closure functional on the 120-vertex 600-cell: Paper~XXII constructs $\Fcal$ as a log-sum-exp over squared distances to the 600-cell vertices,
\begin{equation}
    \Fcal[\Phi] \;=\; -\varepsilon^2 \log \sum_{v \in V_{600}} \exp\!\Bigl(-\frac{\|\Phi - v\|^2}{2\varepsilon^2}\Bigr),
\end{equation}
an $H_4$-invariant soft-minimum over vertex distances (Paper~XXII; parameter $\varepsilon^2$ is distinct from the stochastic noise scale $\sigma^2$ used elsewhere in this paper). In the quantum track, $\Fcal$ drives the stochastic dynamics: the stationary distribution is $P \propto e^{-2\Fcal/\sigma^2}$ (Paper~XIV~\citep{SmartXIV}); Paper~XVIII~\citep{SmartXVIII} pairs forward/backward processes under $\Fcal$-gradient drift via the Nelson construction; and the associated Witten Hamiltonian in its source form is $\HW = -(\sigma^2/2)\Delta + V_{\mathrm{W}}$ with
\begin{equation}
V_{\mathrm{W}} \;=\; \frac{|\nabla \Fcal|^2}{2\sigma^2} - \frac{\Delta \Fcal}{2}
\end{equation}
(Papers XVIII/XIX/XXI). The Witten Hamiltonian governs the \emph{equilibrium-tangent} dynamics only; Paper~XXI notes that \emph{Route A} (Witten Hamiltonian on closure-state space) and \emph{Route B} (graph Laplacian on the 600-cell vertex graph) have different kinetic operators on different spaces and are not presently identified to a single operator.

In the gravitational track, Paper~XXIII~\citep{SmartXXIII} defines admissibility of overlap events \emph{geometrically} from relative birth-angle data, not via closure-residual thresholds on $\Fcal$; any refinement of admissibility by $\Fcal$ is outside the cited construction. The source field of Paper~XXVII~\citep{SmartXXVII} is defined from event density, degree deviation, or geodesic traffic (three explicit source candidates), and the graph-Poisson equation $\Delta_G u = S - \bar{S}$ is introduced as a \emph{model} on the event-order graph, with the mean-subtraction an algebraic centering identity rather than a dynamical conservation law.

\begin{remark}[The two appearances of $\Fcal$ are not identified]
The closure functional $\Fcal$ drives the quantum-track dynamics (Witten potential, Nelson pairing, stationary distribution). In the gravitational track, the cited Papers~XXIII--XXVII do \emph{not} derive admissibility or source content from $\Fcal$: admissibility is geometric (XXIII), source is density/degree/traffic (XXVII). The analogy is a programme-level pattern-correspondence, not a proved identification of the two appearances of $\Fcal$ as the same object.
\end{remark}

\subsection{The transition graph}

The event set $\Ecal$ carries a Hasse-graph structure $\Gcal = (\Ecal, E_H)$ (defined in Paper~XXV~\citep{SmartXXV}, not Paper~XXIII). In the quantum track, Papers~XVIII/XIX use the Route~A Witten Hamiltonian on closure-state space, whose kinetic operator is the closure-state Laplacian $\Delta_A$. Paper~XXI separately posits \emph{Route~B} as a conditional alternative: a self-adjoint generator on the 600-cell vertex graph, involving a different Laplacian $\Delta_B$ and not of Witten-Hamiltonian form. In the gravitational track, Paper~XXVII's field equation uses yet a third Laplacian $\Delta_G$ on the event Hasse graph. The three Laplacians $\{\Delta_A, \Delta_B, \Delta_G\}$ act on three different spaces and are \emph{not} shown to be the same operator in any of the cited source papers.

\begin{definition}[Proposed bookkeeping triple, not a derived common substrate]\label{def:substrate}
The triple $(\Ecal, \Fcal, \Gcal)$ is proposed here as a \emph{bookkeeping labelling} of common names used across the two tracks: $\Ecal$ the event/configuration set (different spaces per track), $\Fcal$ the closure functional (track-specific or absent role: used dynamically in the quantum track, not used to define admissibility/source in the gravitational track), $\Gcal$ the graph carrier (different graphs per track). The cited Papers~XV--XXVII do \emph{not} derive a common substrate; each paper uses the objects in its own track's specific sense. This paper proposes the triple as a candidate organising structure for future identification work; no identification theorem is claimed here.
\end{definition}

%==================================================================
\section{The Quantum Regime}\label{sec:quantum}
%==================================================================

The quantum track is characterised by the following structural features, assembled from the cited source-paper results and interpretive correspondences (each feature's derivational status is in the source paper, not re-derived here).

\subsection{Multiplicity of admissible configurations}

Paper~XV~\citep{SmartXV} constructs a multi-path complexified Onsager--Machlup ansatz on the closure-functional landscape; the result is a \emph{phase-coherent sum-over-paths structure}, not a derivation of quantum superposition. Paper~XV's phase-interference mechanism (relative-phase parameter on path amplitudes) and Paper~XXIII's birth-angle ordering mechanism use different variables and are not bridged in the cited source papers; we treat them as distinct constructions. Paper~XVIII~\citep{SmartXVIII} supplies a \emph{separate} Nelson-pairing construction (forward/backward processes under $\Fcal$-gradient drift) sharing the same equilibrium stationary-amplitude structure; XVIII is not an operator acting on XV's construction.

\subsection{Phase structure and interference}

Paper~XV's path amplitudes are defined by the complexified Onsager--Machlup action; the path-integrated form involves $\Fcal$, and the stationary-distribution amplitude takes the form $A(\psi) = \sqrt{P_{\mathrm{st}}(\psi)} = N \exp(-\Fcal(\psi)/\sigma^2)$ (Paper~XV's definition, from Paper~XIV's stationary distribution $P_{\mathrm{st}} \propto e^{-2\Fcal/\sigma^2}$). This is the amplitude assigned to a stationary-distribution configuration $\psi$, \emph{not} a generic per-path amplitude: paths between the same endpoints carry different Onsager--Machlup actions and interfere in Paper~XV's construction.

\subsection{Unitary-type evolution}

At the Paper~XIV equilibrium, the paired forward/backward process yields the Schr\"odinger equation with Witten Hamiltonian $\HW$ (Paper~XVIII). The evolution is norm-conserving (Paper~XIX~\citep{SmartXIX}). Away from equilibrium, the dynamics is nonlinear, with a modulus-curvature residual $\delta U_{\mathrm{rel}}$ that is \emph{global-phase invariant} (invariant under $\Psi \mapsto e^{i\alpha}\Psi$ for constant $\alpha$; this is modulus-only dependence, not a local gauge invariance) and perturbatively small near equilibrium.

\subsection{Summary}

\begin{quote}
The quantum regime is characterised by: multiplicity of admissible configurations, path-level phase structure producing interference, and norm-conserving evolution governed by the closure functional through the Witten Hamiltonian.
\end{quote}

%==================================================================
\section{The Gravitational Regime}\label{sec:gravitational}
%==================================================================

The gravitational track is characterised by the following structural features, with the conditional status stated in the cited Papers~XXIII--XXVII.

\subsection{Event ordering}

The per-event birth-angle data $\theta_b(e)$ (Paper~XXIII~\citep{SmartXXIII}) defines a strict partial order on admissible events (irreflexive and transitive, basepoint-dependent); non-totality of this order is conditional on the existence of incomparable events in the chosen poset. Under the order, events are ordered by when they become admissible under $\theta_b$; events at the same $\theta_b$-level are incomparable (not ``simultaneously accessible'' in a physical-time sense). Note that $\theta_b$ (Paper XXIII's per-event birth angle) is a different object from Paper~XV's phase variable $\theta$; the two tracks use the same symbol for different quantities.

\subsection{Observer frames and kinematics}

In Paper~XXIV~\citep{SmartXXIV}, an observer is defined as a pair (linear extension of the poset, positive sampling weight function); different observers in this sense define different frames. Frame-dependent simultaneity is exhibited under XXIV's conditional hypothesis that incomparable events exist in the chosen poset. An endpoint-interval time-dilation analogue is exhibited under XXIV's separate conditional hypothesis of unequal weight sums along a shared linear extension. These are \emph{kinematic analogues}, not Lorentz-derived consequences; connection to physical velocity is not supplied in XXIV.

\subsection{Hasse-graph metric, curvature indicators, graph-Poisson source model}

The Hasse graph $\Gcal$ carries a distance structure on connected components (Paper~XXV~\citep{SmartXXV}, not a global Lorentzian metric), graph-curvature indicators on toy models (Paper~XXVI~\citep{SmartXXVI}, reference-dependent and not curvature-tensor-grade), and a graph-Poisson source-potential model $\Delta_G u = S - \bar{S}$ (Paper~XXVII~\citep{SmartXXVII}, introduced as a model definition, not derived from cascade dynamics). The proposed pattern ``source $\to$ potential $\to$ curvature $\to$ geodesic concentration'' is a toy-model observation (Paper~XXVII's bowtie example), not a derived chain: Paper~XXVII explicitly states that its field equation is introduced as a model, that no Einstein equation or stress-energy tensor is derived, and that $u$ does not weight geodesics.

\subsection{Summary}

\begin{quote}
The gravitational regime is characterised by: constraint-driven event ordering, frame-dependent kinematic analogues (under XXIV's conditional hypotheses), Hasse-graph metric on connected components, graph-curvature indicators on toy models, and a source-potential response modelled by the graph-Poisson equation.
\end{quote}

%==================================================================
\section{Structural Correspondence}\label{sec:correspondence}
%==================================================================

The following table records the principal proposed correspondences suggested by the preceding constructions. Each row pairs objects that share a name across the two tracks, with the ``Status'' column recording the strength of the correspondence (formal analogy on different spaces, different mathematical use, etc.). No row records a proved identification of the two tracks' objects as the same mathematical entity.

\begin{table}[ht]
\centering
\caption{Proposed pattern-correspondence between quantum and gravitational tracks. Each row records objects sharing a name across the two tracks; the ``Status'' column flags whether the cited papers identify the object as \emph{the same} mathematical entity (same space, same operator) or only as \emph{formally analogous} (analogous role on different spaces).}
\label{tab:correspondence}
\begin{tabular}{@{}lllc@{}}
\toprule
Shared name & Quantum role & Gravitational role & Status \\
\midrule
Event/configuration set $\Ecal$ & Closure state space / 600-cell vertices & Event set of overlap pairs & Different spaces \\
Closure functional $\Fcal$ & Stochastic potential (XXII form) & Not used for admissibility in XXIII--XXVII & Different use \\
Graph Laplacian $\Delta$ & Route~A Witten kinetic $\Delta_A$ (closure-state) \emph{or} Route~B $\Delta_B$ (600-cell vertex graph, conditional) & Hasse-graph field operator $\Delta_G$ (XXV/XXVII) & Three distinct operators on three spaces \\
Paths & Interference sum over closure paths (XV) & Hasse-graph geodesics (XXV) & Different paths \\
$\Fcal$-gradient vs source $S(e)$ & Nelson drift (XVIII) & Density/degree/traffic (XXVII) & Different objects \\
Norm conservation vs source balance & $\|\Psi\|^2$ const (XIX) & $\sum(S - \bar S) = 0$ algebraic identity & Different laws \\
Equilibrium vs vacuum & Stationary $\Rst^2$ (XIV) & Constant-$u$ configuration (XXVII) & Different states \\
Non-equilibrium vs curvature & Residual $\delta U_{\mathrm{rel}}$ (XIX/XX) & $K_{\mathrm{pot}} > 0$ graph indicator (XXVII) & Different quantities \\
\bottomrule
\end{tabular}
\end{table}

\begin{remark}[Status of the correspondence: pattern-matched, not identified]
Table~\ref{tab:correspondence} is a \emph{pattern-correspondence table}, not an identification theorem. The ``Status'' column records that in every row the quantum-track object and the gravitational-track object live on different spaces or are different mathematical quantities; only their \emph{role} (kinetic operator, potential, path object) is analogous. The cited Papers~XV--XXVII do \emph{not} identify these objects as the same mathematical entity. Whether a genuine identification theorem holds --- for any of the rows --- is an open question for dedicated work, not resolved by this paper.
\end{remark}

\subsection{Path structure: interference vs geodesics}

This correspondence deserves elaboration because it connects the most characteristic features of each regime.

In the quantum track, Paper~XV assigns to each path a complexified Onsager--Machlup action; path amplitudes involve the path-integrated closure functional $\Fcal$, and the stationary-distribution amplitude takes the form $A(\psi) = N \exp(-\Fcal(\psi)/\sigma^2)$ (derived from Paper~XIV's stationary distribution, \emph{not} a per-path amplitude assignment).

In the gravitational track, Paper~XXVII defines graph geodesics as shortest paths on the Hasse graph $\Gcal$ \emph{independently} of the accessibility potential $u(e)$; $u$ does not weight geodesic lengths. The observational toy-model correlation between large $K_{\mathrm{pot}}$ and geodesic concentration is noted in Paper~XXVII's bowtie example but is not derived as a motion law from $u$.

\begin{quote}
Quantum: path amplitudes involve the Onsager--Machlup action and $\Fcal$ (Paper~XV).\\
Gravitational: geodesics are $\Gcal$-shortest-paths, independent of $u$ (Paper~XXVII).
\end{quote}

A candidate interpretation is that a stationary-phase limit of the path-amplitude sum might concentrate on geodesic structure; however, no such limit is derived in Papers~XV or XXVII, and the analogy remains a pattern-suggestion rather than a proved stationary-phase/WKB correspondence.

\subsection{Operator structure: Witten Hamiltonian vs graph Laplacian}

In the quantum track, the equilibrium-tangent Witten Hamiltonian (Papers XVIII/XIX/XXI) is
\begin{equation}
    \HW^{A} \;=\; -\frac{\sigma^2}{2}\,\Delta_A + V_{\mathrm{W}}, \qquad V_{\mathrm{W}} \;=\; \frac{|\nabla\Fcal|^2}{2\sigma^2} - \frac{\Delta_A \Fcal}{2},
\end{equation}
\emph{Route~A}: $\Delta_A$ is the Laplacian on closure-state space (Paper~XVIII), and $\HW^{A}$ above is the Route~A Witten Hamiltonian. \emph{Route~B} (Paper~XXI, conditional): posits a separate self-adjoint generator $H_B$ involving the 600-cell vertex graph Laplacian $\Delta_B$; $H_B$ is \emph{not} a Witten Hamiltonian with $V_{\mathrm{W}}$. Paper~XXI explicitly states that Route~A and Route~B generate dynamics on different spaces and have not been identified as a single operator.

In the gravitational track, the Paper~XXVII field equation is
\begin{equation}
    \Delta_G\, u(e) \;=\; S(e) - \bar{S},
\end{equation}
where $\Gcal$ is the Hasse event-order graph of Paper~XXV~\citep{SmartXXV} and $\Delta_G$ is the graph Laplacian introduced as the field operator in Paper~XXVII~\citep{SmartXXVII}. Paper~XXV defines the Hasse-graph structure and its component-wise metric; the graph-Laplacian/Poisson field equation is a Paper~XXVII construction.

\begin{remark}[The Laplacians are formally analogous, not identified]
The three Laplacians $\Delta_A$ (quantum-track Route~A, closure-state space), $\Delta_B$ (quantum-track Route~B, 600-cell vertex graph), and $\Delta_G$ (gravitational-track Hasse graph) are \emph{formally analogous} operators that act on \emph{three different spaces}. Papers~XV--XXVII do not identify any pair of these as the same operator. In addition to this space mismatch, the three operators are defined with different sign conventions in their source papers (XXVII adopts the combinatorial graph-Laplacian convention $(\Delta_G f)(e) = \sum_{e' \sim e} (f(e') - f(e))$, while Route~B and Route~A follow their own conventions); any prospective identification theorem would therefore have to fix a common convention. The role-analogy --- kinetic term for $\Delta_A$ in the Witten Hamiltonian; separate self-adjoint generator for $\Delta_B$; field-equation operator for $\Delta_G$ --- is a pattern correspondence, not a proved identification.
\end{remark}

%==================================================================
\section{Regime Transition}\label{sec:transition}
%==================================================================

The two regimes are compared here via a finite combinatorial statistic on the event poset: the \textit{ordering fraction} $\eta$. $\eta$ is a discrete-combinatorial indicator, not a derived physical control parameter; whether $\eta$ varies continuously under any physical evolution is an open question not addressed by this paper.

\subsection{The ordering parameter}

\begin{definition}[Constraint ordering strength]\label{def:ordering}
For a finite event set $\Ecal$ with $|\Ecal| \ge 2$ and partial order $\preceq$ (with strict part $\prec$ defined by $e_1 \prec e_2 :\!\iff e_1 \preceq e_2 \wedge e_1 \neq e_2$), the \textit{ordering fraction} is
\begin{equation}
    \eta \;=\; \frac{|\{(e_1, e_2) \in \Ecal \times \Ecal : e_1 \prec e_2\}|}{\binom{|\Ecal|}{2}} \;\in\; [0, 1],
\end{equation}
equivalently, the fraction of unordered event pairs that are comparable under $\preceq$ (each comparable unordered pair contributes exactly one ordered strict pair to the numerator).
\end{definition}

Assuming $|\Ecal| \ge 2$:
\begin{itemize}
    \item $\eta = 0$ (antichain): no comparable pairs. Maximum combinatorial linear-extension count. Candidate indicator for the multiplicity regime with which Papers~XV--XXI's constructions are heuristically associated here; no theorem identifies antichains with ``quantum mechanics'' in any standard sense.
    \item $\eta = 1$ (total order): one linear extension. Candidate indicator for the ordered regime where the order/metric/source parts of the gravitational-track scaffold simplify (single linear extension, well-defined Hasse metric on a connected component, no simultaneity structure from Paper~XXIV since total order makes all pairs comparable); no theorem identifies total orders with ``gravity'' in any standard sense.
    \item $0 < \eta < 1$: intermediate. The quantum-style path-sum structures of XV--XXI and the gravitational-style Hasse-graph structures of XXIII--XXVII can in principle \emph{coexist} as programme labels, but no dynamical crossover between the two is proved.
\end{itemize}

\subsection{Weak-ordering regime: quantum}

This is proposed as the \emph{weak-ordering label} for posets with small $\eta$; no theorem connects small $\eta$ to Paper~XV interference or to Paper~XVIII's Schr\"odinger recovery. Papers~XV--XXI operate on closure-state or 600-cell vertex configurations, not on event posets; the label is a suggestive mapping, not a derivational linkage.

\subsection{Strong-ordering regime: gravitational}

This is proposed as the \emph{strong-ordering label} for posets with $\eta$ near $1$. On such posets, and under the finite-graph connectedness hypothesis of Papers~XXV/XXVII (the graph-distance metric is a metric only on each connected component, and the graph-Poisson compatibility $\sum_e (S(e) - \bar S) = 0$ must be imposed componentwise on each component of the Hasse graph), graph-geodesic language becomes available on the Hasse graph (Paper~XXV), graph-curvature indicators can be defined (Paper~XXVI), and the graph-Poisson source-potential model (Paper~XXVII) can be written down. Dominance of geodesics over other path contributions is \emph{not} proved in Papers~XXIII--XXVII, which operate at the scaffold level; the label is a suggestive mapping, not a derivational linkage.

\subsection{The crossover}

\begin{proposition}[Combinatorial linear-extension count as an $\eta$-indicator]\label{prop:regime}
For a finite poset $(\Ecal, \preceq)$ with $|\Ecal| = n \ge 2$ and ordering fraction $\eta$ (Definition~\ref{def:ordering}), the number $L(\preceq)$ of linear extensions satisfies:
\begin{enumerate}
    \item \textup{(Antichain)} If $\eta = 0$, then $L(\preceq) = n!$ (every permutation is a linear extension).
    \item \textup{(Total order)} If $\eta = 1$, then $L(\preceq) = 1$.
    \item \textup{(Intermediate)} For $0 < \eta < 1$, $L(\preceq)$ lies strictly between $1$ and $n!$, with the specific value determined by the full poset structure (not by $\eta$ alone).
\end{enumerate}
\end{proposition}

\begin{proof}
\emph{(1)} If $\eta = 0$, the strict relation $\prec$ is empty (antichain case; the reflexive order $\preceq$ is still the diagonal, not empty), so every bijection $\Ecal \to \{1, \dots, n\}$ is order-preserving; hence $L(\preceq) = n!$. \emph{(2)} If $\eta = 1$, $\preceq$ is a total order, and the unique order-preserving bijection is the only linear extension; $L(\preceq) = 1$. \emph{(3)} For $0 < \eta < 1$: monotonicity of $L$ under addition of a new \emph{strict} comparable relation (one between previously incomparable elements, not already forced by transitivity) gives $L < n!$ strictly, since each such addition rules out at least one previously-valid linear extension. For the lower bound $L > 1$: pick any incomparable pair $(e, e')$ in $\preceq$; by the order-extension principle, both $e \prec_{\!*} e'$ and $e' \prec_{\!*} e$ can be extended to linear orders compatible with $\preceq$, giving at least two distinct linear extensions. Strict inequalities $1 < L < n!$ therefore hold at any intermediate $\eta$.

\emph{``Not by $\eta$ alone'' counterexample.} Two posets on four elements both with exactly two strict comparable pairs (hence the same $\eta$) can have different linear-extension counts: the poset $P_1 = \{a \prec b,\ c \prec d\}$ (two disjoint chains of length $2$) has $L(P_1) = \binom{4}{2} = 6$ (choose which two of the four positions the $\{a, b\}$-chain occupies, in order). The poset $P_2 = \{a \prec b,\ a \prec c\}$ (``V''-shape: $a$ below both $b$ and $c$, with $d$ independent of all) has $L(P_2) = 8$: $a$ must occupy one of positions $1, 2$ (otherwise $b$ or $c$ would precede $a$); for each position of $a$, count the orderings of $\{b, c, d\}$ in the remaining positions that respect $\preceq$ ($d$ can go anywhere, $b$ and $c$ are both above $a$ with no constraint between them). Enumeration gives $L(P_2) = 8 \neq 6 = L(P_1)$. Thus $L$ depends on the full poset-covering structure, not on $\eta$ alone.
\end{proof}

\begin{remark}[What Proposition~\ref{prop:regime} does and does not establish]
Proposition~\ref{prop:regime} is a purely combinatorial statement about linear-extension counting. It does \emph{not} establish that antichains reproduce quantum mechanics or that total orders reproduce gravity. On the association between $L(\preceq)$ and observers: Paper~XXIV defines an observer as (linear extension, positive sampling weight); $L(\preceq)$ counts linear-extension orderings only. With a fixed canonical weight convention, $L(\preceq)$ counts distinct ordering-frames; if positive weights are chosen freely, there are infinitely many observers per linear extension, so $L(\preceq)$ does \emph{not} bound the observer count in that case. The interpretation of $L(\preceq)$ as a proxy for ``path multiplicity'' or ``quantum'' behaviour is a programme-level label, not a derived correspondence.
\end{remark}

\begin{remark}[Scale and ordering]
The claim is not that ``quantum $=$ small'' and ``gravity $=$ large'' in any absolute sense. The claim is a proposed labelling: as constraint ordering strengthens on a finite event poset (by whatever mechanism; no dynamical mechanism is supplied here), the proposed interpolation moves from the regime in which Papers~XV--XXI's quantum-track constructions are being heuristically associated here with one end of the antichain-to-total-order spectrum, to the regime in which Papers~XXIII--XXVII's gravitational-track constructions are being heuristically associated with the other. The ordering fraction $\eta$ \emph{labels} this proposed interpolation; it does not prove a dynamical transition. Whether $\eta$ correlates with physical scale is an empirical question the framework does not answer.
\end{remark}

%==================================================================
\section{Substrate-Bridge Hypothesis}\label{sec:unification}
%==================================================================

The preceding sections --- the shared bookkeeping triple (Definition~\ref{def:substrate}), the pattern-correspondence table (Table~\ref{tab:correspondence}) of formal analogies across different spaces, and the ordering-fraction label (Definition~\ref{def:ordering}), which is a combinatorial poset statistic and not a derived physical control parameter --- together motivate the following programme-level hypothesis (not established as a theorem by this paper):

\begin{quote}
\textbf{Substrate-bridge hypothesis.} Quantum behaviour (as captured by Papers~XV--XXII's construction) and gravitational structure (as captured by Papers~XXIII--XXVII's scaffold) \emph{may} be viewable as two organisational regimes of a candidate shared event-order geometry whose existence and identification are themselves to be established. The substrate-bridge hypothesis is \emph{not} established by this paper as a theorem; it is proposed as a candidate organising framework subject to identification theorems (operator identification, space identification, $\Fcal$-vs-admissibility identification) that would need to be proved in dedicated follow-up work.
\end{quote}

\begin{remark}[Epistemic status: hypothesis, not theorem]
This is a programme-level hypothesis motivated by the pattern-correspondence table and the ordering-fraction analogy. It is neither a unification proof nor a derivation of quantum gravity. It is not a replacement for either quantum mechanics or general relativity. It identifies a pattern within the VFD programme that would require further mathematical work --- in particular, theorems identifying quantum-track and gravitational-track Laplacians, functionals, and configuration spaces as the same mathematical objects --- to upgrade to a substrate-identification theorem.
\end{remark}

%==================================================================
\section{Crystallisation as a Candidate Bridge}\label{sec:crystallisation}
%==================================================================

The crystallisation operator --- the founding object of the VFD programme --- acquires a new interpretation in this context.

\subsection{The crystallisation functional}

In the original crystallisation formalism (see \texttt{papers/formalism/crystallisation-formalism.tex}), the crystallisation operator $C[\Psi]$ selects from the argmin set:
\begin{equation}
C[\Psi] \;\in\; \operatorname*{arg\,min}_{\Phi \in \mathcal A} F[\Phi].
\end{equation}
The source establishes non-emptiness of this argmin set under finite-dimensional compactness of $\mathcal A$, together with Lyapunov stability of strict local minimisers and monotonic descent along the gradient flow of $F$; the source's infinite-dimensional treatment states compactness/continuity/coercivity more loosely rather than in the sharp weak-compactness-plus-weak-lower-semicontinuity form. The standard mathematical strengthening --- weak-compactness of $\mathcal A$ together with weak-lower-semicontinuity of $F$ --- yields existence by the direct method in infinite dimensions; we adopt this as a \emph{supplementary hypothesis} for precision, beyond the loose compactness/continuity/coercivity statement of the original formalism. Global state-space convergence of the gradient flow is \emph{not} established without additional assumptions. If $\mathcal A$ is additionally taken to be convex and $F$ strictly convex on $\mathcal A$ (a \emph{supplementary hypothesis} not asserted by the original formalism), the minimizer is unique; otherwise $C[\Psi]$ denotes an arbitrary element of the non-empty minimizer set. The functional $F$ in the original formalism is a generic stochastic closure functional (as in Paper~XIV), not necessarily the specific Paper~XXII log-sum-exp $\Fcal$ used on the 600-cell; the use of $\Fcal$ in crystallisation prose throughout this paper is a \emph{name-level} reuse, not a proved identification of the original-formalism $F$ with the XXII functional.

\subsection{Reinterpretation}

Under the substrate-bridge hypothesis (not proved in this paper), crystallisation would be \emph{a candidate description} of such a transition from the quantum-track regime to the gravitational-track regime --- if, and only if, that bridge hypothesis is later established by the identification theorems enumerated in Section~\ref{sec:unification}. In the present paper, the following is \emph{interpretive language conditional on the unproved bridge hypothesis} and \emph{on the absence of any $\eta$-dynamics or map from minimisation of $F$ to strengthening of a poset order}:

\begin{itemize}
    \item \textbf{Before crystallisation (proposed conditional picture)}: multiple admissible configurations coexist. The event set is weakly ordered. Quantum-track path-interference structures are the heuristic label.
    \item \textbf{During crystallisation (proposed conditional picture)}: the original crystallisation functional $F$ would generate gradient-flow descent toward stable minimisers under the formalism's stated hypotheses (monotonic descent and Lyapunov stability; global convergence is not established). Constraint ordering would have to strengthen in any dynamical version of this picture --- $\eta$ would have to increase --- but no cascade dynamics supplies such a mechanism in the cited papers.
    \item \textbf{After crystallisation (proposed conditional picture)}: one would model the post-crystallisation regime as strongly ordered. In that modelled regime, and under the finite-graph connectedness hypothesis of Papers~XXV/XXVII (with the Poisson compatibility applied componentwise on disconnected Hasse graphs), Hasse-graph geodesic, metric, curvature-indicator, and graph-Poisson language becomes available (Papers~XXV--XXVII); dominance and full emergence of continuum field-equation structure are not established.
\end{itemize}

\begin{quote}
Crystallisation is proposed here as a possible geometric bridge by which quantum multiplicity could be modelled as resolving into gravitational order, \emph{conditional on} the substrate-bridge hypothesis being upgraded to a theorem via the identification theorems of Section~\ref{sec:unification}.
\end{quote}

\begin{remark}
This interpretation is structural, not dynamical in the sense of specifying a time-dependent process. The ``before/during/after'' language describes the logical relationship between regimes, not necessarily a temporal sequence. Whether crystallisation occurs in time, or whether time itself emerges from event-order geometry as Paper~XXIII suggests (via birth-angle ordering, not via crystallisation), is a question the framework poses but does not resolve.
\end{remark}

%==================================================================
\section{Discussion and Limitations}\label{sec:discussion}
%==================================================================

\subsection{What this paper establishes and does not establish}

\begin{itemize}
    \item \textbf{Presented}: the pattern-correspondence table (Table~\ref{tab:correspondence}) listing role-parallels between the two tracks' named objects, with the ``Status'' column explicitly flagging each row as formal analogy on different spaces, not mathematical identification.
    \item \textbf{Proposed}: the bookkeeping-triple (Definition~\ref{def:substrate}), labelling $(\Ecal, \Fcal, \Gcal)$ as common names used across the two tracks with track-specific definitions.
    \item \textbf{Established}: Proposition~\ref{prop:regime}, the combinatorial linear-extension count on finite posets (antichain: $n!$ extensions; total order: $1$; intermediate: strictly between) and its $\eta$ labelling.
    \item \textbf{Not established}: any identification theorem between quantum-track and gravitational-track Laplacians, functionals, or configuration spaces. The three Laplacians $\{\Delta_A, \Delta_B, \Delta_G\}$ act on three different spaces (Route~A closure-state, Route~B 600-cell vertex graph, gravitational Hasse graph).
    \item \textbf{Not established}: that $(\Ecal, \Fcal, \Gcal)$ is a \emph{derived} common substrate. The triple is a proposed bookkeeping label.
    \item \textbf{Not established}: that $\eta$ is a derived physical control parameter, or that antichains ``are'' quantum mechanics or total orders ``are'' gravity.
    \item \textbf{Not established}: crystallisation as a proved mechanism for a quantum-to-gravitational transition; the section on crystallisation is interpretive.
\end{itemize}

\subsection{Further open items beyond the above}

\begin{itemize}
    \item No derivation of the Einstein equations from quantum mechanics or vice versa.
    \item No theory of quantum gravity.
    \item No continuum limit. Whether the discrete structures converge to continuous quantum field theory or general relativity is open.
    \item No experimental predictions from the substrate-bridge hypothesis itself (though both constituent tracks make independent structural predictions, documented in their source papers).
    \item No proof that the pattern-correspondence of Table~\ref{tab:correspondence} is unique or necessary; alternative role-mappings may exist.
    \item $\eta$ is defined on finite posets; continuum/thermodynamic-limit behaviour is not analysed.
    \item No sharp phase-transition threshold between multiplicity and ordering regimes.
\end{itemize}

\subsection{Comparison with existing approaches}

\begin{table}[ht]
\centering
\caption{Informal shorthand comparison with other programmes. The external-approach rows (string theory, loop quantum gravity, causal set theory) summarise programme-origin characterisations without primary citations; they are \emph{not} intended as evaluative status assessments of those programmes and would require dedicated reviews with primary sources to license. Only the VFD row is supported by the cited papers of this programme.}
\label{tab:comparison}
\begin{tabular}{@{}llll@{}}
\toprule
Approach & Starting point & Unification mechanism & Source basis in this paper \\
\midrule
String theory & Continuous strings & Extra dimensions, dualities & Informal shorthand \\
Loop quantum gravity & Spin networks & Discrete spacetime & Informal shorthand \\
Causal set theory & Discrete partial orders & Order $\to$ geometry & Informal shorthand \\
VFD (this paper) & Closure on 600-cell & Substrate-bridge hypothesis & Hypothesis-level (see Sec.~\ref{sec:unification}) \\
\bottomrule
\end{tabular}
\end{table}

\begin{remark}
Among the approaches listed in Table~\ref{tab:comparison}, VFD is closest in spirit to causal set theory, which also begins from discrete partial orders. The key structural difference is that VFD constructs a basepoint-dependent strict order from overlap geometry on the dual 600-cell (Paper~XXIII), whereas causal set theory takes the partial order as fundamental. Whether this difference is physically significant is an open question. The comparison to string theory and loop quantum gravity in Table~\ref{tab:comparison} is at the level of informal programme-origin shorthand; detailed comparison to those frameworks is outside this paper's scope and would require dedicated treatment with primary citations.
\end{remark}

%==================================================================
\section{Programme Summary}\label{sec:programme}
%==================================================================

The VFD closure programme's current documented structure is as follows:

\begin{table}[ht]
\centering
\caption{Current documented programme structure. Entries characterise the \emph{currently documented} scope of each track, not an external status assessment. Paper~XXII in particular presents structural/numerical correspondences around the 600-cell closure functional, not a derivation of mass, gauge couplings, or constants ``from'' $\Fcal$ in any fundamental-derivation sense.}
\label{tab:programme}
\begin{tabular}{@{}llll@{}}
\toprule
Track & Papers & Core documented content & Regime \\
\midrule
Crystallisation & Flagship + A--D & Deterministic state-selection formalism & Selection \\
Mass / constraint & I--XIII & Particle-mass correspondences with 600-cell spectral data & Spectral \\
Quantum recovery & XV--XXI & Schr\"odinger recovery on the equilibrium tangent and conditional Route~B & Quantum \\
Spectral bridge & XXII & Mass/gauge/constants correspondences around the 600-cell $\Fcal$ & Spectral--quantum \\
GR scaffold & XXIII--XXVII & Event order, Hasse metric, graph-curvature indicators, graph-Poisson source model & Gravitational scaffold \\
Unification & XXVIII (this paper) & Structural bridge \emph{hypothesis} across regimes & Both \\
\bottomrule
\end{tabular}
\end{table}

%==================================================================
\section{Future Directions}\label{sec:future}
%==================================================================

The structural bridge hypothesised here opens several directions for future work:

\begin{enumerate}
    \item \textbf{Continuum limit.} Does the discrete framework converge to known continuum theories (quantum field theory, general relativity) in appropriate limits? This is the defining technical challenge.

    \item \textbf{Regime-transition dynamics.} The ordering fraction $\eta$ labels the proposed crossover picture, but no dynamics for $\eta$ is supplied. A dynamical theory --- specifying how $\eta$ would evolve under a physical process and how this evolution would connect quantum-track and gravitational-track constructions --- would substantially strengthen the framework.

    \item \textbf{Coupling constant.} The gravitational field equation $\Delta_G u = S - \bar{S}$ has no coupling constant (proportionality set to 1). Identifying a Newton-constant analogue from the 600-cell spectral data is an open target.

    \item \textbf{Observable signatures.} Both tracks make structural predictions independently. Whether the unification principle itself generates new, testable predictions --- beyond those of the constituent tracks --- is the key question for experimental relevance.
\end{enumerate}

%==================================================================
\section{Conclusion}\label{sec:conclusion}
%==================================================================

The VFD closure programme has constructed two related tracks: a quantum track (Papers~XV--XXII) in which Schr\"odinger recovery is exact at Paper~XIV equilibrium and formal at leading order near equilibrium (Papers XVIII/XIX/XXI), and a gravitational track (Papers~XXIII--XXVII) supplying discrete scaffold analogues of temporal ordering, observer kinematics, Hasse-graph metric separation, graph-curvature indicators, and a graph-Poisson source-potential model.

This paper argues, at \emph{structural-correspondence level}, that these two tracks can be jointly viewed under a shared set of name-level objects (bookkeeping triple $(\Ecal, \Fcal, \Gcal)$ with track-specific definitions). Table~\ref{tab:correspondence} records role-parallels between the two tracks; the ``Status'' column flags every row as a formal analogy on different spaces, not a mathematical identification. Proposition~\ref{prop:regime} proves only combinatorial linear-extension-count bounds, not a physical regime-transition theorem.

The result is a \emph{substrate-bridge hypothesis}, not a theory of quantum gravity. Upgrading the pattern to a substrate-identification theorem requires dedicated mathematical work: operator-identification among $\{\Delta_A, \Delta_B, \Delta_G\}$ (three distinct Laplacians currently acting on three different spaces, with no pair identified in the cited source papers), space-identification between closure-state/600-cell and the event Hasse graph, and a proved link between the quantum-track $\Fcal$ and the gravitational-track admissibility/source content. These are open items of the next stage of the programme.

What can be stated is this: the distinction between quantum and gravitational structure in the VFD framework is proposed here to be organisable by event accessibility; whether it is fundamentally reducible to event accessibility remains an open question for the identification theorems listed above.

%==================================================================
\bibliographystyle{plainnat}
\bibliography{references}

\end{document}
